\subsection{Compact-Density State and Transition Events}

For each frame, the swarm center and root-mean-square radius were computed as
\begin{equation}
R_{\mathrm{rms}}(t)=
\sqrt{
\frac{1}{n(t)}
\sum_i \left\|\mathbf{x}_i(t)-\bar{\mathbf{x}}(t)\right\|^2
}.
\end{equation}
We then used the radius-based density proxy
\begin{equation}
\rho_{\mathrm{rms}}(t)=
\frac{n(t)}{(4/3)\pi R_{\mathrm{rms}}(t)^3}.
\end{equation}
Within each observation, $\rho_{\mathrm{rms}}(t)$ was robust-z standardized to
define the compact-density coordinate $C(t)$. A one-second centered rolling
mean of this standardized coordinate provided the smoothed compact-density
trace, and $\dot{C}(t)$ was computed as the time gradient of that smoothed
trace. The swarm-size coordinate $R(t)$ was the robust-z standardized
$R_{\mathrm{rms}}(t)$.

Transition events were inherited from the upstream compact-density
coarse-graining. The continuous quantity was $C(t)$ and the low/high labels
were the frozen \texttt{spectral\_set} state labels. A transition was defined
as a switch between the two labels. It was retained only when both the
preceding and following state runs lasted at least 0.20 s, a persistence
screen used to remove single-frame or very short label flicker at the 100 Hz
sampling rate. The event time was the first frame of the new run, and events
were labeled as low-to-high or high-to-low transitions.
The \texttt{spectral\_set} labels came from an upstream transfer-operator
compact-density coarse-graining of robust-standardized \texttt{r\_rms},
\texttt{density\_rms}, and \texttt{anisotropy}. The selected \texttt{eig2}
partition separated a more compact high-density state from a less compact
low-density state and was constructed upstream of T1, not fitted from T1.

\subsection{Local Non-Affine T1 Observable}

T1 is the frozen local tangential non-affine measurement target used in this
study. It is not a raw individual velocity, an inferred force, or a behavioral
rule. It is computed from focal-neighborhood relative motion after local
affine deformation has been removed.

For a focal individual $i$, let
\begin{equation}
\mathbf{r}_{ij}(t)=\mathbf{x}_j(t)-\mathbf{x}_i(t), \qquad
\Delta\mathbf{r}_{ij}(t,\ell)
=
\mathbf{r}_{ij}(t+\ell)-\mathbf{r}_{ij}(t),
\end{equation}
where $j$ indexes retained nearest neighbors of the focal individual and
$\ell$ is the finite temporal lag. The local affine deformation was fitted by
equal-weight least squares,
\begin{equation}
\hat{\mathbf{J}}_i(t)=
\arg\min_{\mathbf{J}}
\sum_{j\in N_i(t)}
\left\|
\Delta\mathbf{r}_{ij}(t,\ell)
-
\mathbf{J}\mathbf{r}_{ij}(t)
\right\|^2 .
\end{equation}
Because the fit uses focal-centered relative coordinates, local translation is
already removed. The finite-lag non-affine residual velocity is
\begin{equation}
\mathbf{u}^{NA}_{ij}(t)=
\frac{
\Delta\mathbf{r}_{ij}(t,\ell)
-
\hat{\mathbf{J}}_i(t)\mathbf{r}_{ij}(t)
}{\Delta t_{\ell}},
\end{equation}
where $\Delta t_{\ell}$ is the elapsed time corresponding to lag $\ell$.

Let $\hat{\boldsymbol{\rho}}_i(t)$ denote the unit vector from the swarm center
to focal individual $i$. The tangential residual is
\begin{equation}
\mathbf{u}^{NA}_{ij,\mathrm{tan}}(t)=
\mathbf{u}^{NA}_{ij}(t)
-
\left[
\mathbf{u}^{NA}_{ij}(t)\cdot\hat{\boldsymbol{\rho}}_i(t)
\right]
\hat{\boldsymbol{\rho}}_i(t).
\end{equation}
The primary positive analyses used unsigned tangential magnitude or activity.
For the later state-matched analyses, the focal-centered activity was frozen
as
\begin{equation}
A^{NA}_{\mathrm{tan},i}(t)=
\operatorname{median}_{j\in N_i(t)}
\left\|
\mathbf{u}^{NA}_{ij,\mathrm{tan}}(t)
\right\|^2,
\end{equation}
and the swarm-level frame activity was
\begin{equation}
A^{NA}_{\mathrm{tan}}(t)=
\operatorname{median}_{i} A^{NA}_{\mathrm{tan},i}(t).
\end{equation}
We refer to this residual activity family as T1. Signed quantities were used
only for secondary event-type and moment-structure tests.
